\chapter{Дүгнэлт болон Цаашдын хөгжүүлэлт}
\section{Ажлын дүгнэлт}
Энэхүү бакалаврын судалгааны ажлын хүрээнд Sumbee.mn сошиал мониторингийн системд тулгарч буй хэрэгцээ шаардлага болох агуулгын шинжилгээний нарийвчлалыг сайжруулах зорилгоор Topic Modeling болон NER (Named Entity Recognition)-ийн нэгтгэсэн шийдлийг боловсрууллаа. 

Боловсруулсан загварыг инференс пайплайнаар ажиллуулж туршсан үр дүн нь энгийн мэдлэгийн сан (Knowledge Database) үүсгэхээс илүү контекст таних чадвар сайтай, LLM API ашигласнаас хэд дахин хямд өртөгтэй бөгөөд хурдтай гэдгийг батлав. Уг шийдэл нь сошиал тренд болон үүсч буй чухал сэдвүүдийг автоматжуулсан байдлаар гаргахад чухал ач холбогдолтой юм.

\sectionsection{Цаашид хөгжүүлэх боломж}
Судалгааны ажлыг цаашид дараах чиглэлээр сайжруулах боломжтой:
\begin{itemize}
    \item Олон хэлний загвар (Multilingual) болон монгол хэл дээр тусгайлан fine-tune хийгдсэн Topic Classification загварыг аппликэйшнд бодит хугацааны инференст (Real-time inference) оруулах.
    \item Сургагдсан моделуудыг API хэлбэрт оруулан Sumbee.mn-ийн микросервис архитектурт (Microservice) бүрэн нэгтгэх.
    \item Network graph-д temporal (цаг хугацааны) анализ нэмж, чиг хандлагын өөрчлөлтийг урьдчилан таамагладаг болгох.
\end{itemize}
